\begin{bt}%[1D7N2-1]
% nguon: Archive/2024/TOAN_11_XE/TOAN_11_XE/data4/11C9/KNTT-9.2.tex #1295

	Tính đạo hàm của hàm số $y=x^5$ tại điểm $x=2$ và $x=-\dfrac{1}{2}$.
	\loigiai{Ta có $\left(x^5\right)^\prime =5 x^4$. Từ đó, $y^\prime \left(2\right)=5\cdot 2^4=80$ và $y^\prime \left(-\dfrac{1}{2}\right)=5\cdot \left(-\dfrac{1}{2}\right) ^4=\dfrac{5}{16}$.
	}
\end{bt}

\begin{bt}%[1D7N2-1]
% nguon: Archive/2024/TOAN_11_XE/TOAN_11_XE/data4/11C9/KNTT-9.2.tex #1296

	Tính đạo hàm của hàm số $y=x^{10}$ tại điểm $x=-1$ và $x=\sqrt[3]{2}$.
	\loigiai{Ta có $\left(x^{10}\right)^\prime =10 x^9$. Từ đó, $y^\prime \left(-1\right)=10\cdot \left(-1\right)^9=-10$ và $y^\prime \left(\sqrt[3]{2}\right)=10\cdot \left(\sqrt[3]{2}\right) ^9=80$.
	}
\end{bt}

\begin{bt}%[1D7N2-1]
% nguon: Archive/2024/TOAN_11_XE/TOAN_11_XE/data4/11C9/KNTT-9.2.tex #1297

	Tính đạo hàm của hàm số $y=\sqrt{x}$  tại điểm $x=1$ và $x=\dfrac{1}{4}$.
	\loigiai{Ta có $y^\prime=\left(\sqrt{x}\right)^\prime =\dfrac{1}{2\sqrt{x}}$, $x>0$. Từ đó, $y^\prime \left(1\right)=\dfrac{1}{2\sqrt{1}}=\dfrac{1}{2}$ và $y^\prime \left(\dfrac{1}{4}\right)=\dfrac{1}{2\sqrt{\dfrac{1}{4}}}=\dfrac{1}{2\cdot \dfrac{1}{2}}=1$.
	}
\end{bt}

\begin{bt}%[1D7N2-3]
% nguon: Archive/2024/TOAN_11_XE/TOAN_11_XE/data4/11C9/KNTT-9.2.tex #1298

	Viết phương trình tiếp tuyến của đồ thị hàm số $y=\sqrt{x}$ tại điểm có hoành độ bằng $4$.
	\loigiai{Với $x=4$ thì $y=\sqrt{4}=2$.\\
		Ta có $y^\prime=\left(\sqrt{x}\right)^\prime =\dfrac{1}{2\sqrt{x}}$, $x>0$. Từ đó, $y^\prime \left(4\right)=\dfrac{1}{2\sqrt{4}}=\dfrac{1}{4}$. \\
		Phương trình tiếp tuyến của đồ thị hàm số $y=\sqrt{x}$ tại điểm có hoành độ bằng $4$ là 
		\[y-2=\dfrac{1}{4}\left(x-4\right) \ \text{hay }\ y=\dfrac{1}{4}x+1.\]	
	} 
\end{bt}

\begin{bt}%[1D7N2-1]
% nguon: Archive/2024/TOAN_11_XE/TOAN_11_XE/data4/11C9/KNTT-9.2.tex #1299

	Tìm đạo hàm của hàm số $y=\sqrt[3]{x}$ tại điểm $x=8$.
	\loigiai{
		Ta có $y'=\left(\sqrt[3]{x}\right)'=\left(x^{\frac{1}{3}}\right)^\prime=\dfrac{1}{3}\cdot x^{\frac{1}{3}-1}=\dfrac{1}{3}\cdot x^{-\frac{2}{3}}=\dfrac{1}{3\sqrt[3]{x^2}}$.\\
		Từ đó, $y'(8)=\dfrac{1}{3\sqrt[3]{8^2}}=\dfrac{1}{12}$.
	}
\end{bt}

\begin{bt}%[1D7N2-1]
% nguon: Archive/2024/TOAN_11_XE/TOAN_11_XE/data4/11C9/KNTT-9.2.tex #1300

	Tìm đạo hàm của các hàm số
	\begin{listEX}[2]
		\item $y=\sqrt[4]{x}$ tại $x=1$;
		\item $y=\dfrac{1}{x}$ tại $x=-\dfrac{1}{4}$.
	\end{listEX}
	\loigiai{
		\begin{listEX}
			\item $y=\sqrt[4]{x}$ tại $x=1$.\\
			
			Ta có $y^\prime =\left(\sqrt[4]{x}\right)^\prime=\left(x^{\frac{1}{4}}\right)^\prime=\dfrac{1}{4}\cdot x^ {\frac{1}{4}-1}=\dfrac{1}{4}\cdot x^ {-\frac{3}{4}}=\dfrac{1}{4\sqrt[4]{x^3}}$.
			Từ đó, $y^\prime(1)=\dfrac{1}{4\sqrt[4]{1^3}}=\dfrac{1}{4}$.
			\item $y=\dfrac{1}{x}$ tại $x=-\dfrac{1}{4}$.\\
			Ta có $y^\prime =\left(\dfrac{1}{x}\right)^\prime=-\dfrac{1}{x^2}$. Từ đó, $y^ \prime \left(-\dfrac{1}{4}\right)=-\dfrac{1}{\left(-\frac{1}{4}\right)^2}=-16$.
		\end{listEX}
	}
\end{bt}

\begin{bt}%[1D7N2-1]
% nguon: Archive/2024/TOAN_11_XE/TOAN_11_XE/data4/11C9/KNTT-9.2.tex #1301

	Tính đạo hàm của hàm số $y=\cos x$ tại $x=\dfrac{\pi}{6}$.
	\loigiai{Ta có $y^\prime= \left(\cos x\right)^\prime =-\sin x$. Vậy $y^\prime\left(\dfrac{\pi}{6}\right)=-\sin \dfrac{\pi}{6}=-\dfrac{1}{2}$.
	}
\end{bt}

\begin{bt}%[1D7N2-1]
% nguon: Archive/2024/TOAN_11_XE/TOAN_11_XE/data4/11C9/KNTT-9.2.tex #1302

	Tính đạo hàm của hàm số $y=\tan x$ tại $x=\dfrac{3\pi}{4}$.
	\loigiai{Ta có $y^\prime= \left(\tan x\right)^\prime =\dfrac{1}{\cos^2x}\ \left(x \neq \dfrac{\pi}{2}+k\pi, k \in \mathbb{Z}\right)$. Vậy $y^\prime\left(\dfrac{3\pi}{4}\right)= \dfrac{1}{\cos^2 \left(\frac{3\pi}{4}\right)}=2$.
	}
\end{bt}

\begin{bt}%[1D7N2-1]
% nguon: Archive/2024/TOAN_11_XE/TOAN_11_XE/data4/11C9/KNTT-9.2.tex #1303

	Tìm đạo hàm của các hàm số 
	\begin{listEX}[2]
		\item $y=\mathrm{e}^x$ tại $x=2\ln 3$;
		\item $y=\log_5x $ tại $x=2$.
	\end{listEX}
	\loigiai{Tìm đạo hàm của các hàm số 
		\begin{listEX}[1]
			\item $y=\mathrm{e}^x$ tại $x=2\ln 3$.\\
			Ta có $y'=\left(\mathrm{e}^x\right)^\prime =\mathrm{e}^x$. Từ đó, $y^\prime\left(2\ln 3\right)=\mathrm{e}^{2\ln 3}=\left(\mathrm{e}^{\ln 3}\right)^2=3^2=9$.
			\item $y=\log_5x $ tại $x=2$.\\
			Ta có $y'=\left( \log_5 x \right)^\prime =\dfrac{1}{x \cdot \ln 5 } \  \left(x>0\right)$.
			Từ đó $y'(2)=\dfrac{1}{2 \cdot \ln 5 }$.
	\end{listEX}}
\end{bt}

\begin{bt}%[1D7N2-1]
% nguon: Archive/2024/TOAN_11_XE/TOAN_11_XE/data4/11C9/KNTT-9.2.tex #1304

	Tìm đạo hàm của các hàm số 
	\begin{listEX}[2]
		\item $y=9^x$ tại $x=1$;
		\item $y=\ln x$ tại $x=\dfrac{1}{3}$.
	\end{listEX}
	\loigiai{Tìm đạo hàm của các hàm số 
		\begin{listEX}
			\item $y=9^x$ tại $x=1$.\\
			Ta có $y'=\left(9^ x\right)^\prime =9^x \cdot \ln 9$. Từ đó $y'(1)=9^1\cdot \ln 9=18\ln 3$.
			\item $y=\ln x$ tại $x=\dfrac{1}{3}$.\\
			Ta có $y'=\left(\ln x\right)^\prime =\dfrac{1}{x}\ \left(x>0\right)$.  Từ đó $y'\left(\dfrac{1}{3}\right)=\dfrac{1}{\frac{1}{3}}=3$.
		\end{listEX}
	}
\end{bt}

\begin{bt}%[1D7H2-1]
% nguon: Archive/2024/TOAN_11_XE/TOAN_11_XE/data4/11C9/KNTT-9.2.tex #1305

	Tính đạo hàm của các hàm số sau:
	\begin{listEX}[3]
		\item $y=3x^2-4x+2$;
		\item $y=x\sin x$;
		\item $y=\dfrac{3x+2}{2x-1}$.
	\end{listEX}
	\loigiai{
		\begin{listEX}
			\item $y=3x^2-4x+2$.\\
			$y'=\left(3x^2-4x+2\right)^\prime=\left(3x^2\right)^\prime-\left(4x\right)^\prime+(2)^\prime=3\left(x^2\right)^\prime-4\left(x\right)^\prime+0=3\cdot2x-4\cdot 1=6x-4$.
			\item $y=x\sin x$.\\
			$y'=\left(x\sin x\right)^\prime=\left(x\right)'\sin x+x \left(\sin x\right)'=1\cdot \sin x+x \cdot \cos x=\sin x+x\cos x$.
			\item $y=\dfrac{3x+2}{2x-1}$.
			\begin{eqnarray*}
				y'&=&\left(\dfrac{3x+2}{2x-1}\right)^\prime=\dfrac{\left(3x+2\right)'\left(2x-1\right)-\left(3x+2\right)\left(2x-1\right)'}{\left(2x-1\right)^2}\\
				&=&\dfrac{3\cdot \left(2x-1\right)-(3x+2)\cdot 2}{\left(2x-1\right)^2}
				=\dfrac{6x-3-6x-4}{\left(2x-1\right)^2}=\dfrac{-7}{\left(2x-1\right)^2}.
			\end{eqnarray*}
	\end{listEX}}
\end{bt}

\begin{bt}%[1D7H2-1]
% nguon: Archive/2024/TOAN_11_XE/TOAN_11_XE/data4/11C9/KNTT-9.2.tex #1306

	Tính đạo hàm của các hàm số sau:
	\begin{listEX}[2]
		\item $y=x^2\cdot 3^x$;
		\item $y=\dfrac{\sqrt{x}}{\cos x}$.
	\end{listEX}
	\loigiai{\begin{listEX}
			\item $y=x^2\cdot 3^x$.\\
			$y'=\left(x^2\cdot 3^x\right)'=\left(x^2\right)'\cdot 3^x+x^2\cdot \left(3^x\right)'=2x\cdot 3^x+x^2\cdot 3^x\cdot \ln 3=x\cdot 3^x\left(2+x\ln 3\right)$.
			\item $y=\dfrac{\sqrt{x}}{\cos x}$.
			\begin{eqnarray*}
				y'&=&\left(\dfrac{\sqrt{x}}{\cos x}\right)'=\dfrac{\left(\sqrt{x}\right)'\cdot \cos x-\sqrt{x}\cdot \left(\cos x\right)'}{\cos^2x}=\dfrac{\dfrac{1}{2\sqrt{x}}\cdot\cos x-\sqrt{x}\left(-\sin x \right) }{\cos^2x}\\
				&=&\dfrac{\cos x+2x\sin x}{2\sqrt{x}\cos^2x}.
			\end{eqnarray*}
			
		\end{listEX}
	}
\end{bt}

\begin{bt}%[1D7H2-1]
% nguon: Archive/2024/TOAN_11_XE/TOAN_11_XE/data4/11C9/KNTT-9.2.tex #1307

	Tính đạo hàm của các hàm số sau:
	\begin{listEX}[2]
		\item $y=x\cdot \log_2x$;
		\item  $y=x^3\cdot \mathrm{e}^x$.
	\end{listEX}	
	\loigiai{\begin{listEX}
			\item $y=x\cdot \log_2x$.\\
			$y'=\left(x\cdot \log_2x\right)'=x'\cdot \log_2x+x\cdot \left(\log_2x\right)'=1\cdot \log_2x+x \cdot \dfrac{1}{x\ln 2}=\log_2x+\dfrac{1}{\ln 2}$.
			\item  $y=x^3\cdot \mathrm{e}^x$.\\
			$y'=\left(x^3\cdot \mathrm{e}^x\right)'=\left(x^3\right)'\cdot \mathrm{e}^x+x^3\cdot \left(\mathrm{e}^x\right)'=3x^2\cdot \mathrm{e}^x+x^3\cdot \mathrm{e}^x=x^2\mathrm{e}^x\left(3+x\right)$.
		\end{listEX}
	}
\end{bt}

\begin{bt}%[1D7H2-1]
% nguon: Archive/2024/TOAN_11_XE/TOAN_11_XE/data4/11C9/KNTT-9.2.tex #1308

	Tính đạo hàm của các hàm số sau:
	\begin{listEX}[3]
		\item $y=(3x^2+x)^3$;
		\item $y=\sin 2x$;
		\item $y=\mathrm{e}^{x^2+1}$.
		
	\end{listEX}
	\loigiai{\begin{enumerate}
			\item $y=(3x^2+x)^3$.\\
			Đặt $u=3x^2+x$ thì $y=u^3$. Ta có $u'_x=6x+1$ và $y'_u=3u^2$.\\
			Suy ra $y'_x=y'_u\cdot u'_x=3u^2\cdot (6x+1)=3\left(3x^2+x\right)^2(6x+1)$.\\
			Vậy $y'=3\left(3x^2+x\right)^2(6x+1)$.
			\item $y=\sin 2x$.\\
			Đặt $u=2x$ thì $y=\sin u$. Ta có $u'_x=2$ và $y'_u=\cos u$.\\
			Suy ra $y'_x=y'_u\cdot u'_x=\cos u\cdot 2=\cos 2x\cdot 2 =2\cos 2x$.\\
			Vậy $y'=2\cos 2x$.
			\item $y=\mathrm{e}^{x^2+1}$.\\
			Đặt $u=x^2+1$ thì $y=\mathrm{e}^u$. Ta có $u'_x=2x$ và $y'_u
			=\mathrm{e}^u$.\\
			Suy ra $y'_x=y'_u\cdot u'_x=\mathrm{e}^u\cdot 2x=\mathrm{e}^{x^2+1} \cdot 2x =2x\mathrm{e}^{x^2+1}$.\\
			Vậy $y'=2x\mathrm{e}^{x^2+1}$.
	\end{enumerate}}
\end{bt}

\begin{bt}%[1D7H2-1]
% nguon: Archive/2024/TOAN_11_XE/TOAN_11_XE/data4/11C9/KNTT-9.2.tex #1309

	Tính đạo hàm của các hàm số sau:
	\begin{listEX}[3]
		\item $y=\left(2x^3+3\right)^2$;
		\item $y=\cos 3x$;
		\item $y=\log_2\left(x^2+2\right)$.
	\end{listEX}
	\loigiai{
		\begin{enumerate}
			\item $y=\left(2x^3+3\right)^2$.\\
			Đặt $u=2x^3+3$ thì $y=u^2$. Ta có $u'_x=6x$ và $y'_u
			=2u$.\\
			Suy ra $y'_x=y'_u\cdot u'_x=2u\cdot 6x=2\left(2x^3+3\right)\cdot 6x =12x\left(2x^3+3\right)$.\\
			Vậy $y'=12x\left(2x^3+3\right)$.
			\item $y=\cos 3x$.\\
			Đặt $u=3x$ thì $y=\cos u$. Ta có $u'_x=3$ và $y'_u
			=-\sin u$.\\
			Suy ra $y'_x=y'_u\cdot u'_x=\left(-\sin u\right)\cdot 6=\left(-\sin 3x\right)\cdot 3 =-3\sin 3x$.\\
			Vậy $y'=-3\sin 3x$.
			\item $y=\log_2\left(x^2+2\right)$.\\
			Đặt $u=x^2+2$ thì $y=\log_2 u$. Ta có $u'_x=2x$ và $y'_u
			=\dfrac{1}{u\ln 2}$.\\
			Suy ra $y'_x=y'_u\cdot u'_x=\dfrac{1}{u\ln 2} \cdot 2x=\dfrac{1}{\left(x^2+2\right)\ln2}\cdot 2x=\dfrac{2x}{\left(x^2+2\right)\ln 2}$.\\
			Vậy $y'=\dfrac{2x}{\left(x^2+2\right)\ln 2}$.
	\end{enumerate}}
\end{bt}

\begin{bt}%[1D7N3-1]
% nguon: Archive/2024/TOAN_11_XE/TOAN_11_XE/data4/11C9/KNTT-9.3.tex #1310

	Tính đạo hàm cấp hai của các hàm số:
	\begin{multicols}{2}
		\begin{enumerate}
			\item $y=3x^2+5x+1$;
			\item $y=\sin x$.
		\end{enumerate}
	\end{multicols}
	\loigiai{
		\begin{enumerate}
			\item $y=3x^2+5x+1$.\\
			$y'=\left(3x^2+5x+1\right)'=6x+5$, $y^{\prime\prime} =\left(6x+5\right)'=6$.
			\item $y=\sin x$.\\
			$y'=\left(\sin x\right)'=\cos x$, $y^{\prime\prime} =\left(\cos x\right)'=-\sin x$.
		\end{enumerate}
	}
\end{bt}

\begin{bt}%[1D7H3-1]
% nguon: Archive/2024/TOAN_11_XE/TOAN_11_XE/data4/11C9/KNTT-9.3.tex #1311

	Tính đạo hàm cấp hai của các hàm số sau
	\begin{listEX}[2]
		\item $y=2x^4-5x^2+3$;
		\item $y=x\cdot \mathrm{e}^x$.
	\end{listEX}
	\loigiai{\begin{enumerate}
			\item $y=2x^4-5x^2+3$.\\
			Ta có $y'=\left(2x^4-5x^2+3\right)'=8x^3-10x$; $y^{\prime\prime}=\left(8x^3-10x\right)'=24x^2-10$.
			\item $y=x\cdot \mathrm{e}^x$.\\
			Ta có $y'=\left(x\cdot \mathrm{e}^x\right)'=x'\cdot \mathrm{e}^x+x \cdot \left(\mathrm{e}^x\right)'=1\cdot \mathrm{e}^x+x\cdot \mathrm{e}^x=(1+x)\cdot \mathrm{e}^x$.\\
			$y^{\prime\prime}=\left[(1+x)\cdot \mathrm{e}^x\right]'=(1+x)'\cdot \mathrm{e}^x+(1+x)\cdot \left(\mathrm{e}^x\right)'=\mathrm{e}^x+(1+x)\mathrm{e}^x=(2+x)\mathrm{e}^x$.
		\end{enumerate}
	}
\end{bt}

\begin{bt}%[1D7H3-1]
% nguon: Archive/2024/TOAN_11_XE/TOAN_11_XE/data4/11C9/KNTT-9.3.tex #1313

	Tính đạo hàm cấp hai của các hàm số sau:
	\begin{multicols}{2}
		\begin{enumerate}
			\item $y=x^2-x$;
			\item $y=\cos x$.
		\end{enumerate}
	\end{multicols}	
	\loigiai{\begin{enumerate}
			\item $y=x^2-x$.\\
			$y'=(x^2-x)'=2x-1$, $y^{\prime\prime}=(2x-1)'=2$.
			\item $y=\cos x$.\\
			$y'=(\cos x)'=-\sin x$, $y^{\prime\prime}=(-\sin x)'=-(\sin x)'=-\cos x$.
		\end{enumerate}
	}
\end{bt}

\begin{bt}%[1D7H1-4]
% nguon: Archive/2024/TOAN_11_XE/TOAN_11_XE/data4/11C9/KNTT-9.3.tex #1314

	Một hòn sỏi rơi tự do có quãng đường rơi tính theo thời gian $t$ là $s(t)=4,9 t^2$, trong đó $s$ tính bằng mét và $t$ tính bằng giây. Tính gia tốc rơi của hòn sỏi lúc $t=3$.
	\loigiai{
		Ta có $s'(t)=\left(4{,}9 t^2\right)'=4{,}9\cdot 2t=9{,}8t$;  $s^{\prime\prime}(t)=(9{,}8t)'=9{,}8$.\\
		Gia tốc rơi của hòn sỏi lúc $t=3$ là $s^{\prime\prime}(3)=9{,}8$ m/s$^2$.
	}
\end{bt}

\begin{bt}%[1D7H2-8]
% nguon: Archive/2024/TOAN_11_XE/TOAN_11_XE/data4/11C9/KNTT-9.3.tex #1315

	Một công ty xác định rằng tổng chi phí của họ, tính theo nghìn đô-la, để sản xuất $x$ mặt hàng là $C(x)=\sqrt{5x^2+60}$ và công ty lên kế hoạch nâng sản lượng trong $t$ tháng kể từ nay theo hàm số $x(t)=20t+40$. Chi phí sẽ tăng thế nào sau $4$ tháng kể từ khi công ty thực hiện kế hoạch đó?
	\loigiai{Ta có $C'(x)=\dfrac{(5x^2+60)'}{2\sqrt{2x^2=60}}=\dfrac{5x}{\sqrt{5x^2+50}}$.\\
		Công ty thực hiện kế hoạch nâng sản lượng  sau $4$ tháng  có sản lượng là $x(4)=20\cdot 4+40=120$.\\
		Chi phí sẽ tăng sau $4$ tháng kể từ khi công ty thực hiện kế hoạch là $C'(120)\simeq 2,235$.}
\end{bt}

\begin{bt}%[1D7H1-4]
% nguon: Archive/2024/TOAN_11_XE/TOAN_11_XE/data4/11C9/KNTT-9.3.tex #1316

	Trên Mặt Trăng, quãng đường rơi tự đo của một vật được cho bởi công thức $s(t)=0,81t^2$, trong đó $t$ là thời gian được tính bằng giây và $s$ tính bằng mét. Một vật thả rơi từ độ cao $200~m$  phía trên Mặt Trăng. Tại thời điểm $t=2$ sau khi thả vật đó, tính:
	\begin{multicols}{2}
		\begin{enumerate}
			\item Quãng đường vật đã rơi;
			\item Gia tốc của vật.
		\end{enumerate}
	\end{multicols}	
	\loigiai{Ta có $s'(t)=1,62t$, $s''(t)=1,62$.
		\begin{enumerate}
			\item Quãng đường vật đã rơi tại thời điểm $t=2$ sau khi thả vật rơi từ độ cao $200$ m  phía trên Mặt Trăng là $s(2)=0,81\cdot 2^2=3,24$ m.
			\item Gia tốc của vật thời điểm $t=2$ là $s''(2)=1,62$ m/s$^2$.
		\end{enumerate}
	}
\end{bt}

\begin{bt}%[1D7H1-3]
% nguon: Archive/2024/TOAN_11_XE/TOAN_11_XE/data4/11C9/KNTT-9.1.tex #1328

	Cho hàm số $y=f(x)=x^2$ có đồ thị $(C)$ và điểm $M(2 ; 4) \in(C)$. Tính hệ số góc của tiếp tuyến của $(C)$ tại điểm $M$ và viết phương trình của tiếp tuyến đó.
	\loigiai{
		Ta có $\left(x^2\right)^{\prime}=2 x$ nên tiếp tuyến của $(C)$ tại $M$ có hệ số góc là $f^{\prime}(2)=2 \cdot 2=4$.
		Phương trình tiếp tuyến của $(C)$ tại $M$ là
		$$
		y-4=4(x-2) \Leftrightarrow y=4 x-4 .
		$$}
\end{bt}

\begin{bt}%[1D7H1-3]
% nguon: Archive/2024/TOAN_11_XE/TOAN_11_XE/data4/11C9/KNTT-9.1.tex #1329

	Cho $(C)$ là đồ thị của hàm số $f(x)=\dfrac{1}{x}$ và điểm $M(1 ; 1) \in(C)$. Tính hệ số góc của tiếp tuyến của $(C)$ tại điểm $M$ và viết phương trình tiếp tuyến đó.
	\loigiai{
		Đạo hàm $f^\prime(x)=-\dfrac{1}{x^2}$.
		\begin{itemize}
			\item Hệ số góc của tiếp tuyến $(C)$ tại $M(1 ; 1)$ là $k=f^\prime(1)=-1.$
			\item Phương trình tiếp tuyến của $(C)$ tại $M(1 ; 1)$ là $y=-x+2.$
	\end{itemize}}
\end{bt}

\begin{bt}%[1D7H1-1]
% nguon: Archive/2024/TOAN_11_XE/TOAN_11_XE/data4/11C9/KNTT-9.1.tex #1330

	Dùng định nghĩa để tính đạo hàm của các hàm số sau
	\begin{listEX}[3]
		\item $f(x)=-x^2$;
		\item $f(x)=x^3-2 x$;
		\item $f(x)=\dfrac{4}{x}$.
	\end{listEX}
	\loigiai{Tính đạo hàm của các hàm số sau
		\begin{enumerate}
			\item Với bất kì $x_0$, ta có\\
			$f^\prime(x_0)=\lim\limits_{x\to x_0}\dfrac{f(x)-f(x_0)}{x-x_0}=\lim\limits_{x\to x_0}\dfrac{-x^2+{x_0}^2}{x-x_0}=\lim\limits_{x\to x_0}\left[-(x+x_0)\right]=-2x_0$.\\
			Vậy $f^\prime(x)=-2x$ trên $\mathbb{R}$.
			\item Với bất kì $x_0$, ta có\\
			$f^\prime(x_0)=\lim\limits_{x\to x_0}\dfrac{f(x)-f(x_0)}{x-x_0}=\lim\limits_{x\to x_0}\dfrac{\left[\left(x^3-2x\right)-\left({x_0}^3-2x_0\right)\right]}{x-x_0}$\\
			$=\lim\limits_{x\to x_0}\left[x^2+xx_0+{x_0}^2-2\right]=3{x_0}^2-2$.\\
			Vậy $f^\prime(x)=3x^2-2$ trên $\mathbb{R}$.
			\item Với mọi $x_0\ne 0$, ta có\\
			$f^{\prime}\left(x_0\right)=\lim\limits_{x \to x_0} \dfrac{\dfrac{4}{x}-\dfrac{4}{x_0}}{x-x_0}=\lim\limits_{x \to x_0} \dfrac{4(x_0-x)}{x x_0\left(x-x_0\right)}=\lim\limits_{x \to x_0} \dfrac{-4}{x x_0}=-\dfrac{4}{x_0^2}.$\\
			Vậy $f^{\prime}(x)=\left(\dfrac{4}{x}\right)^{\prime}=-\dfrac{4}{x^2}$ trên các khoảng $(-\infty ; 0)$ và $(0 ;+\infty)$.
	\end{enumerate}}
\end{bt}

\begin{bt}%[1D7H1-3]
% nguon: Archive/2024/TOAN_11_XE/TOAN_11_XE/data4/11C9/KNTT-9.1.tex #1331

	Cho hàm số $f(x)=-2 x^2$ có đồ thị $(C)$ và điểm $A(1 ;-2) \in(C)$. Tính hệ số góc của tiếp tuyến với $(C)$ tại điểm $A$.
	\loigiai{Đạo hàm $f^\prime(x)=-4x$.\\
		Hệ số góc tiếp tuyến của $(C)$ tại điểm $A$ là
		$$k=f^\prime(1)=-4\cdot 1=-4.$$
	}
\end{bt}

\begin{bt}%[1D7H2-3]
% nguon: Archive/2024/TOAN_11_XE/TOAN_11_XE/data4/11C9/KNTT-9.1.tex #1332

	Viết phương trình tiếp tuyến của đồ thị hàm số $y=x^3$.
	\begin{multicols}{2}
	\begin{enumerate}
		\item Tại điểm $(-1 ; 1)$;
		\item Tại điểm có hoành độ bằng $2$.
	\end{enumerate}
	\end{multicols}	
	\loigiai{
		\begin{enumerate}
			\item Ta có $f^\prime(x)=3x^2$ nên $f^\prime(-1)=3$.\\
			Phương trình tiếp tuyến của đồ thị tại điểm $(-1;1)$ là
			$$y=f^\prime(-1)\cdot (x+1)+1\Leftrightarrow y=3x+4.$$
			\item Gọi $A(2;8)$ thuộc đồ thị hàm số $y=x^3$\\
			Ta có $f^\prime(x)=3x^2$ nên $f^\prime(2)=12$.\\
			Phương trình tiếp tuyến của đồ thị tại $A$ là
			$$y=f^\prime(2)\cdot (x-2)+8\Leftrightarrow y=12x-16.$$
	\end{enumerate}}
\end{bt}

\begin{bt}%[1D7H1-4]
% nguon: Archive/2024/TOAN_11_XE/TOAN_11_XE/data4/11C9/KNTT-9.1.tex #1333

	Một chuyển động thẳng xác định bởi phương trình $s(t)=4 t^3+6 t+2$, trong đó $s$ tính bằng mét và $t$ là thời gian tính bằng giây. Tính vận tốc tức thời của chuyển động tại $t=2$.
	\loigiai{
		Ta có: $v(t)=s^\prime(t)=12t^2+6$.\\
		Vận tốc tức thời tại $t=2$ là $v(2)=12\cdot 2^2+6=54 \;m/s$.
	}	
\end{bt}

\begin{bt}%[1D7H1-5]
% nguon: Archive/2024/TOAN_11_XE/TOAN_11_XE/data4/11C9/KNTT-9.1.tex #1334

	Một người gửi tiết kiệm khoản tiền $10$ triệu đồng vào một ngân hàng với lãi suất $5 \% /$ năm. Tính tổng số tiền vốn và lãi mà người đó nhận được sau một năm, nếu tiền lãi được tính theo thể thức
	\begin{multicols}{2}
		\begin{enumerate}
			\item Lãi kép với kì hạn $6$ tháng;
			\item Lãi kép liên tục.
		\end{enumerate}
	\end{multicols}	
	\loigiai{Số tiền nhận được sau 
		\begin{enumerate}
			\item Lãi kép với kì hạn $6$ tháng.\\
			Số tiền nhận được sau $6$ tháng đầu là 
			$T=10000000\cdot \mathrm{e}^{0.05\cdot \frac{1}{2}}=10253151{,}21$ (đồng).\\
			Số tiền nhận được sau $6$ tháng tiếp theo là 
			$T=10253151,21\cdot \mathrm{e}^{0.05\cdot \frac{1}{2}}=10512710{,}97$ (đồng).
			\item Lãi kép với kì hạn liên tục là $T=10000000\cdot \mathrm{e}^{0.05}=10512710{,}96$ (đồng).
	\end{enumerate}}
\end{bt}

\begin{bt}%[1D7N2-3]
% nguon: Archive/2024/TOAN_11_XE/TOAN_11_XE/data4/11C9/KNTT-9.OTC.tex #1335
%[Tăng Xuân Phú]
	Cho hàm số $f(x)=x^2-2 x+3$ có đồ thị $(C)$ và điểm $M(-1 ; 6) \in(C)$. Viết phương trình tiếp tuyến với $(C)$ tại điểm $M$.
	\loigiai{Ta có $f^\prime(x)=2x-2$ nên $f^\prime(-1)=-4$.\\
		Phương trình tiếp tuyến của đồ thị $(C)$ tại điểm $M(-1 ; 6)$ là
		$$y=f^\prime(-1)(x+1)+6\Leftrightarrow y=-4x+2.$$}
\end{bt}

\begin{bt}%[1D7H2-1]
% nguon: Archive/2024/TOAN_11_XE/TOAN_11_XE/data4/11C9/KNTT-9.OTC.tex #1336
%[Tăng Xuân Phú]
	Tính đạo hàm của các hàm số sau
	\begin{multicols}{3}
	\begin{enumerate}
		\item $y=3 x^4-7 x^3+3 x^2+1$
		\item $y=\left(x^2-x\right)^3$
		\item $y=\dfrac{4 x-1}{2 x+1}$.
	\end{enumerate}
	\end{multicols}	
	\loigiai{Đạo hàm của các hàm số sau là
		\begin{enumerate}
			\item $y^\prime=12x^3-21x^2+6x$.
			\item $y^\prime=3(x^2-x)^2(x^2-x)^\prime=3(x^2-x)^2(2x-1)$.
			\item $y^\prime=\dfrac{(4x-1)^\prime\cdot(2x+1)-(2x+1)^\prime\cdot(4x-1)}{(2x+1)^2}=\dfrac{6}{(2x+1)^2}$.
	\end{enumerate}}
\end{bt}

\begin{bt}%[1D7H2-1]
% nguon: Archive/2024/TOAN_11_XE/TOAN_11_XE/data4/11C9/KNTT-9.OTC.tex #1337
%[Tăng Xuân Phú]
	Tính đạo hàm của các hàm số sau
	\begin{multicols}{2}
	\begin{enumerate}
		\item $y=\left(x^2+3 x-1\right) e^x$;
		\item  $y=x^3 \log _2 x$.
	\end{enumerate}
	\end{multicols}	
	\loigiai{Đạo hàm của các hàm số sau là
		\begin{enumerate}
			\item $y^\prime=(2x+3)e^x+(x^2+3x-1)e^x=e^x(x^2+5x-1)$.
			\item $y^\prime=(x^3)^\prime\ \cdot\log_2 x+x^3\cdot (\log_2 x)^\prime=3x^2\log_2 x+\dfrac{x^3}{x\ln2}$.
	\end{enumerate}}
\end{bt}

\begin{bt}%[1D7H2-1]
% nguon: Archive/2024/TOAN_11_XE/TOAN_11_XE/data4/11C9/KNTT-9.OTC.tex #1338
%[Tăng Xuân Phú]
	Tính đạo hàm của các hàm số sau
	\begin{multicols}{3}
		\begin{enumerate}
			\item $y=\tan \left(e^x+1\right)$;
			\item $y=\sqrt{\sin 3 x}$;
			\item $y=\cot \left(1-2^x\right)$.
		\end{enumerate}
	\end{multicols}	
	\loigiai{Đạo hàm của các hàm số sau là
		\begin{enumerate}
			\item $y^\prime=\dfrac{(e^x+1)^\prime}{\cos^2(e^x+1)}=\dfrac{e^x}{\cos^2(e^x+1}$.
			\item $y^\prime=\dfrac{(\sin3x)^\prime}{2\sqrt{\sin3x}}=\dfrac{3\cos3x}{2\sqrt{\sin3x}}$.
			\item $y^\prime=\dfrac{(1-2^x)^\prime}{\sin^2(1-2^x)}=\dfrac{2^x\ln x}{\sin^2(1-2^x)}$.
	\end{enumerate}}
\end{bt}

\begin{bt}%[1D7H3-1]
% nguon: Archive/2024/TOAN_11_XE/TOAN_11_XE/data4/11C9/KNTT-9.OTC.tex #1339
%[Tăng Xuân Phú]
	Tính đạo hàm cấp hai của các hàm số sau
	\begin{multicols}{2}
	\begin{enumerate}
		\item $y=x^3-4 x^2+2 x-3$;
		\item $y=x^2 e^x$.
	\end{enumerate}
	\end{multicols}	
	\loigiai{Đạo hàm cấp $1$ của các hàm số sau là
		\begin{enumerate}
			\item $y^\prime=3x^2-8x+2$.
			\item $y^\prime=2xe^x+x^2e^x=e^x(2x+x^2)$.
		\end{enumerate}
		Đạo hàm cấp $2$ của các hàm số trên là
		\begin{enumerate}
			\item $y^{\prime\prime}=6x-8$.
			\item $y^{\prime\prime}=e^x(2x+x^2)+e^x(2+2x)=e^x(x^2+4x+2)$.
	\end{enumerate}}
\end{bt}

\begin{bt}%[1D7H1-4]
% nguon: Archive/2024/TOAN_11_XE/TOAN_11_XE/data4/11C9/KNTT-9.OTC.tex #1340
%[Tổng Nguyễn]
	Một viên sỏi rơi từ độ cao $44,1$ m thì quãng đường rơi được biểu diễn bởi công thức $s(t)=4,9t^2$, trong đó $t$ là thời gian tính bằng giây và $s$ tính bằng mét. Tính
	\begin{multicols}{2}
	\begin{enumerate}
		\item Vận tốc rơi của viên sỏi lúc $t=2$;
		\item  Vận tốc của viên sỏi khi chạm đất.
	\end{enumerate}
	\end{multicols}	
	\loigiai{Ta có $s'(t)=4,9\cdot 2t$.
		\begin{enumerate}
			\item Vận tốc rơi của viên sỏi lúc $t=2$ là $s'(2)=4,9\cdot 2\cdot 2=19,6$ m/s.
			\item  Vận tốc của viên sỏi khi chạm đất.\\
			Thời gian viên sỏi rơi từ độ cao $44,1$ m đến lúc chạm đất
			\[4,9t^2=44,1\ \text{suy ra}\ t=3 \ (\text{vì} \ t>0). \]
			Vận tốc của viên sỏi khi chạm đất là $s'(3)=4,9\cdot 2\cdot 3=29,4$ m/s.
		\end{enumerate}
	}
\end{bt}

\begin{bt}%[1D7H1-4]
% nguon: Archive/2024/TOAN_11_XE/TOAN_11_XE/data4/11C9/KNTT-9.OTC.tex #1341
%[Tổng Nguyễn]
	Một vật chuyển động trên một  đường thẳng được xác định bởi công thức $s(t)=2t^3+4t+1$, trong đó $t$ là thời gian tính bằng giây và $s$ tính bằng mét. Tính vận tốc và gia tốc của vật khi $t=1$.
	\loigiai{ Ta có $s'(t)=6t^2+4$, $s''(t)=12t$.\\
		Vận tốc của vật khi $t=1$ là $s'(1)=6\cdot 1^2+4=10$ m/s.\\
		Gia tốc của vật khi $t=1$ là $s''(1)=12\cdot 1=12 ~ m/s^2$.
	}
\end{bt}

\begin{bt}%[1D7H2-8]
% nguon: Archive/2024/TOAN_11_XE/TOAN_11_XE/data4/11C9/KNTT-9.OTC.tex #1342
%[Tổng Nguyễn]
	Dân số $P$ (tính theo nghìn người) của một thành phố nhỏ được tính theo công thức $P(t)=\dfrac{500t}{t^2+9}$, trong đó $t$ là thời gian tính bằng năm. Tìm tốc độ tăng dân số tại thời điểm $t=12$.
	\loigiai{
		$P'(t)=\left(\dfrac{500t}{t^2+9}\right)'=\dfrac{(500t)'(t^2+9)-500t\left(t^2+9\right)'}{\left(t^2+9\right)^2}=\dfrac{500(t^2+9)-100t\cdot 2t}{\left(t^2+9\right)^2}=\dfrac{-500t^2+4500}{(t^2+9)^2}$.\\
		Tốc độ tăng dân số tại thời điểm $t=12$ là $P'(12)\simeq -2,88$ nghìn người/năm.
	}
\end{bt}

\begin{bt}%[1D7H2-8]
% nguon: Archive/2024/TOAN_11_XE/TOAN_11_XE/data4/11C9/KNTT-9.OTC.tex #1343
%[Tổng Nguyễn]
	Hàm số $S(r)=\dfrac{1}{r^4}$ có thể được sử dụng để xác định sức cản $S$ của dòng máu có bán kính $r$ (tính theo milimét) (theo Bách khoa toàn thư Y học "Harrison's internal medicine  $21$st edition"). Tìm tốc độ thay đổi của $S$ theo $r$ khi $r=0,8$.
	\loigiai{
		Ta có $S'(r)=\left(\dfrac{1}{r^4}\right)'=\left(r^{-4}\right)'=-4r^{-5}=\dfrac{-4}{r^5}$.\\
		Tìm tốc độ thay đổi của $S$ theo $r$ khi $r=0,8$ là $S'(0,8)\simeq-12,207$.
	}
\end{bt}

\begin{bt}%[1D7H2-8]
% nguon: Archive/2024/TOAN_11_XE/TOAN_11_XE/data4/11C9/KNTT-9.OTC.tex #1344
%[Tổng Nguyễn]
	Nhiệt độ cơ thể của một người trong thời gian bi bệnh được cho bởi công thức $T(t)=-0,1t^2+1,2t+98,6$, trong đó $T$ là nhiệt độ (tính theo đơn vị đo nhiệt độ Fahrenheit) tại thời điểm $t$ (tính theo ngày). Tìm tốc độ thay đổi của nhiệt độ ở thời điểm $t=1,5$.
	\loigiai{ Ta có $T'(t)=-0,2t+1,2$.\\
		Tốc độ thay đổi của nhiệt độ ở thời điểm $t=1,5$ là $T'(1,5)=0,9$.
	}
\end{bt}

\begin{bt}%[1D7H2-8]
% nguon: Archive/2024/TOAN_11_XE/TOAN_11_XE/data4/11C9/KNTT-9.OTC.tex #1345
%[Tổng Nguyễn]
	Hàm số $R(v)=\dfrac{600}{v}$ có thể dùng để xác định nhịp tim $R$ của một người mà tim của người đó có thể lấy đi được $6000~ml$ máu trên mỗi phút và $v~ml$ máu trên mỗi nhịp đập (theo Bách khoa toàn thư Y học "Harrison's internal medicine  $21$st edition"). Tìm tốc độ thay đổi của nhịp tim khi lượng máu tim đẩy đi ở một nhịp là $v=80$.
	\loigiai{Ta có $R'(v)=-\dfrac{600}{v^2}$.\\
		Tốc độ thay đổi của nhịp tim khi lượng máu tim đẩy đi ở một nhịp $v=80$ là $R'(80)=-0,09375$. 
	}
\end{bt}
